\chapter*{\MakeUppercase{Các từ viết tắt}}
% \addcontentsline{toc}{chapter}{{\bf \MakeUppercase{List of Abbreviations}}}

\begin{longtable}{p{0.15\textwidth} p{0.75\textwidth}}
% \caption{List of Abbreviations} \label{tab:abbreviations} \\
\toprule
\textbf{Từ viết tắt} & \textbf{Mô tả} \\
\midrule
GUI & Giao diện đồ họa người dùng (Graphical User Interface) \\
CNF & Dạng chuẩn hội (Conjunctive Normal Form) \\
SAT & Bài toán tính thỏa mãn của công thức logic \\
UNSAT & Công thức không có phép gán thỏa mãn \\
DPLL & Davis--Putnam--Logemann--Loveland \\
JSON & Định dạng trao đổi dữ liệu JavaScript Object Notation \\
API & Giao diện lập trình ứng dụng \\
BFS & Tìm kiếm theo chiều rộng (Breadth-First Search) \\
\bottomrule
\end{longtable}



\chapter*{\MakeUppercase{Các ký hiệu toán học}}
% \addcontentsline{toc}{chapter}{\MakeUppercase{List of Symbols}}


\begin{longtable}{p{0.2\textwidth} p{0.7\textwidth}}
% \caption{List of Mathematical Notations and Abbreviations} \label{tab:math_abbreviations} \\
\toprule
\textbf{Kí hiệu} & \textbf{Mô tả} \\
\midrule
\(N\) & Kích thước cạnh bàn cờ; $N\in\{3,4,5\}$ \\
\(C_i\) & Mệnh đề nhân vật tại ô $i$ là tội phạm \\
\(KB_t\) & Cơ sở tri thức công khai ở bước $t$ \\
\(R_t\) & Tập ô đã mở ở bước $t$ \\
\(V_t\) & Tập trạng thái đã biết ở bước $t$ \\
\(L_t\) & Tập manh mối công khai ở bước $t$ \\
\(\models\) & Quan hệ hệ quả logic \\
\(M\) & Số lượng biến mệnh đề phụ được cấp phát; $M \in \mathbb{Z}_{\ge 0}$ \\
\(A_j\) & Biến mệnh đề phụ thứ $j$; $j \in \{1, 2, \dots, M\}$ \\
\(\Phi_t\) & Công thức logic mệnh đề dạng CNF biểu diễn $KB_t$ ở bước $t$ \\
\(\ell(i, \sigma)\) & Hàm ánh xạ chỉ số ô $i$ và trạng thái $\sigma$ thành nghĩa văn số nguyên \\
\(\Omega\) & Bộ chỉ số các ô thuộc vùng tham chiếu của một manh mối \\
\(\Lambda\) & Dãy nghĩa văn đầu vào của ràng buộc đếm số lượng \\
\(s_{r, j}\) & Biến phụ của bộ đếm tuần tự tại phần tử thứ $r$ và ngưỡng $j$ \\
\(g_r\) & Biến phụ xác định ô gốc duy nhất trong ràng buộc liên thông \\
\(p_{r, d}\) & Biến phụ biểu thị ô thứ $r$ đạt tới được sau tối đa $d$ bước kề \\
\(\mathcal{N}_\Omega(r)\) & Tập chỉ số các ô trong $\Omega$ kề bốn hướng vuông góc với ô thứ $r$ \\
\(\top\) & Hằng đúng logic (tương đương tập công thức CNF rỗng $\emptyset$) \\
\(\bot\) & Hằng sai logic (tương đương mệnh đề rỗng $()$, trạng thái UNSAT) \\
\(V\) & Tổng số biến đã khai báo trong công thức CNF; $V = N^2 + M$ \\
\(J(\ell)\) & Điểm trọng số Jeroslow--Wang của nghĩa văn $\ell$ trên các mệnh đề chưa \texttt{SAT} \\
\bottomrule
\end{longtable}


\newpage
